\subsection{成周、洛邑与一处向东的政治支点}

在关中取得胜利，并不等于能在东方稳定地统治。商的旧都和旧有政治关系主要分布在更东面的区域；若王室每逢危机都从宗周临时出兵，路途、补给和信息传递都会把胜利消耗掉。因而，\eventref{WZ-1040-DONGZHENG}{西周·三监东征}之后营建成周，首先是一项把军事实力变成可持续在场的安排：王室在洛阳盆地一带取得一个能会集贵族、安置人群、举行礼仪并向东方调兵的节点。

传世文献常把这座城称作“洛邑”或“成周”。两种名称在不同文本中的用法并不总能一一对应；后世所谓王城、成周的空间关系，又叠加了东周以后的城市变化。把它们一概画成一座边界清晰、功能固定的“新首都”，会把尚不能解决的考古和文献问题提前解决。较稳妥的理解是：宗周仍是王室根基，成周则使王权在东方拥有一处更近的礼仪、军事和行政支点；这是双重重心，而不是现代国家意义上的迁都。

这一支点的作用并非只在城墙之内。道路、粮食、役夫、贵族随从、青铜工匠和祭祀所需物资，必须持续向此汇集。城的建立因而把关中、洛阳盆地和更东面的封国接成一条资源链。它也使王室能够更频繁地把册命、赏赐和征调送往东方，却同时增加了运输、守卫和协调的日常开支。

\subsection{把殷遗民迁入新秩序}

传统叙事说周人在平定东方后迁置“殷民”或“殷顽民”，并把这项措施同成周联系起来。这个说法的重要性不在于把人群想成可以整齐移动的棋子，而在于它显示征服者面对的现实：旧商的贵族、工匠、祭祀传统和地方关系不可能随着一场战役自动消失。将其中一部分人纳入新的城邑和新的政治关系，既可能削弱旧有动员网络，也可能保留周初统治所需的技术、劳力与地方知识。

但“迁殷民”不能被写成一项已知人数、路线和完成日期的工程。传世材料多从周王室的命令和正当性出发，考古材料则主要显现墓葬、器物和作坊等可保存的痕迹；它们都难以逐户说明谁被迁走、谁留下、谁以何种身份加入新秩序。所谓“殷遗民”本身也不必然是同质群体，其中不同家族、职业和地方共同体的处境很可能不同。

从制度效果看，迁置是一种重新编组关系的尝试，而不仅是人口移动。被安置者要面对新的服役、居住和祭祀秩序；接收地的居民也要承担供应、土地调整和治安压力。对王室而言，迁置提供了掌握旧商人群和技术资源的可能；对被迁置者及其周边社会而言，它意味着原有亲族与生计网络可能被打断。材料不足以量化其痛苦和规模，却不应因此把这笔成本从叙事中删去。

\begin{sourcenote}[“迁殷民”能证实到什么程度]
\sourcebook{尚书·召诰}、\sourcebook{尚书·洛诰}及相关周初篇章保存了王室经营东方、告诫和安置的政治语言；《史记》的完整叙事成书更晚。洛阳一带西周遗存与墓地说明此地确有新的精英活动和物质投入，但墓葬形制或器物风格不能自动等同于某个人的“族属”。将文献、金文和考古并读，可以确认东方重组是真问题；至于迁置的规模、具体名单和每个群体的经历，仍须保留空白。
\end{sourcenote}
